% sec5_2.tex — Section 5.2: The Fixed-Effects Estimator

\section{The Fixed-Effects Estimator}
\label{sec:5.2}

\subsection*{The Formula}

As already mentioned, the fixed-effects estimator is defined for the transformed
system. Multiplying both sides of \eqref{eq:5.1.1pp} on page~329 from the left by $\mathbf{Q}$ and
making use of the fact that $\mathbf{Q}\mathbf{1}_M = \mathbf{0}$, we obtain
\[
  \mathbf{Q}\mathbf{y}_i = \mathbf{Q}\mathbf{F}_i \boldsymbol{\beta} + \mathbf{Q}\boldsymbol{\eta}_i
\]
or
\begin{equation}
  \tilde{\mathbf{y}}_i = \widetilde{\mathbf{F}}_i \boldsymbol{\beta} + \tilde{\boldsymbol{\eta}}_i
  \quad (i = 1, 2, \ldots, n),
  \label{eq:5.2.1} \tag{5.2.1}
\end{equation}
where
\begin{equation}
  \underset{(M \times 1)}{\tilde{\mathbf{y}}_i} \equiv \mathbf{Q}\mathbf{y}_i, \quad
  \underset{(M \times \#\boldsymbol{\beta})}{\widetilde{\mathbf{F}}_i} \equiv \mathbf{Q}\mathbf{F}_i, \quad
  \underset{(M \times 1)}{\tilde{\boldsymbol{\eta}}_i} \equiv \mathbf{Q}\boldsymbol{\eta}_i.
  \label{eq:5.2.2} \tag{5.2.2}
\end{equation}

Thus, the fixed effect $\alpha_i$ and common regressors drop out in the transformed equations.
Now form the pooled sample of transformed $\mathbf{y}_i$ and $\mathbf{F}_i$ as
\begin{equation}
  \underset{(Mn \times 1)}{\tilde{\mathbf{y}}}
  \equiv \begin{bmatrix} \tilde{\mathbf{y}}_1 \\ \vdots \\ \tilde{\mathbf{y}}_n \end{bmatrix}, \quad
  \underset{(Mn \times \#\boldsymbol{\beta})}{\widetilde{\mathbf{F}}}
  \equiv \begin{bmatrix} \widetilde{\mathbf{F}}_1 \\ \vdots \\ \widetilde{\mathbf{F}}_n \end{bmatrix}.
  \label{eq:5.2.3} \tag{5.2.3}
\end{equation}

The \textbf{fixed-effects estimator} of $\boldsymbol{\beta}$, denoted $\hat{\boldsymbol{\beta}}_{\text{FE}}$,
is the pooled OLS estimator, that
is, the OLS estimator applied to the pooled sample $(\tilde{\mathbf{y}}, \widetilde{\mathbf{F}})$
of size $Mn$. Therefore,
\begin{align}
  \hat{\boldsymbol{\beta}}_{\text{FE}}
  &\equiv (\widetilde{\mathbf{F}}'\widetilde{\mathbf{F}})^{-1}\widetilde{\mathbf{F}}'\tilde{\mathbf{y}}
  \notag \\
  &= \biggl(\frac{1}{n}\sum_{i=1}^{n}\widetilde{\mathbf{F}}_i'\widetilde{\mathbf{F}}_i\biggr)^{\!-1}
     \frac{1}{n}\sum_{i=1}^{n}\widetilde{\mathbf{F}}_i'\tilde{\mathbf{y}}_i
  \notag \\
  &= \biggl[\frac{1}{n}\sum_{i=1}^{n}(\mathbf{Q}\mathbf{F}_i)'(\mathbf{Q}\mathbf{F}_i)\biggr]^{\!-1}
     \frac{1}{n}\sum_{i=1}^{n}(\mathbf{Q}\mathbf{F}_i)'(\mathbf{Q}\mathbf{y}_i)
  \notag \\
  &= \biggl(\frac{1}{n}\sum_{i=1}^{n}\mathbf{F}_i'\mathbf{Q}\mathbf{Q}\mathbf{F}_i\biggr)^{\!-1}
     \frac{1}{n}\sum_{i=1}^{n}\mathbf{F}_i'\mathbf{Q}\mathbf{Q}\mathbf{y}_i
  \notag \\
  &= \biggl(\frac{1}{n}\sum_{i=1}^{n}\mathbf{F}_i'\mathbf{Q}\mathbf{F}_i\biggr)^{\!-1}
     \frac{1}{n}\sum_{i=1}^{n}\mathbf{F}_i'\mathbf{Q}\mathbf{y}_i
  \quad \text{(since $\mathbf{Q}$ is symmetric and idempotent).}
  \label{eq:5.2.4} \tag{5.2.4}
\end{align}

Substituting \eqref{eq:5.2.1} into \eqref{eq:5.2.4}, the sampling error can be obtained as
\begin{align}
  \hat{\boldsymbol{\beta}}_{\text{FE}} - \boldsymbol{\beta}
  &= \biggl(\frac{1}{n}\sum_{i=1}^{n}\widetilde{\mathbf{F}}_i'\widetilde{\mathbf{F}}_i\biggr)^{\!-1}
     \frac{1}{n}\sum_{i=1}^{n}\widetilde{\mathbf{F}}_i'\tilde{\boldsymbol{\eta}}_i
  \notag \\
  &= \biggl(\frac{1}{n}\sum\mathbf{F}_i'\mathbf{Q}\mathbf{F}_i\biggr)^{\!-1}
     \frac{1}{n}\sum\mathbf{F}_i'\mathbf{Q}\boldsymbol{\eta}_i.
  \label{eq:5.2.5} \tag{5.2.5}
\end{align}

Because it is based on deviations from group means, the fixed-effects estimator
is also known as the \textbf{within estimator} or the \textbf{covariance estimator}. Another,
and perhaps more popular, derivation is to apply OLS to levels, i.e., not in deviations
but with group ($i$) specific dummies added to the list of regressors. For this
reason the estimator is also called the \textbf{least squares dummy variables (LSDV)}
estimator. The LSDV interpretation allows us to see what additional assumptions
are needed for the error-components model in developing a finite-sample theory
for the estimator. This is left as Analytical Exercise~1.

\subsection*{Large-Sample Properties}

As you will prove in Analytical Exercise~2, the fixed-effects estimator can also be
written as a GMM estimator, which lends itself to a straightforward development
of large-sample theory for the fixed-effects estimator summarized in Proposition
5.1 below. Proving the proposition, however, can be accomplished more easily by
inspecting the expression \eqref{eq:5.2.5} for sampling error.

\begin{proposition}[large-sample properties of the fixed-effects estimator]
\label{prop:5.1}
Consider the error-components model (consisting of \eqref{eq:5.1.1pp} on page~329,
\eqref{eq:5.1.2}, \eqref{eq:5.1.8}, \eqref{eq:5.1.4}--\eqref{eq:5.1.6}, and
\eqref{eq:5.1.15}), but relax the SUR assumption \eqref{eq:5.1.8b} by
requiring only that
\begin{equation}
  \E(\mathbf{f}_{im} \cdot \eta_{ih}) = \mathbf{0} \quad
  \textit{for all } m, h \;(= 1, 2, \ldots, M),
  \label{eq:5.2.6} \tag{5.2.6}
\end{equation}
where $\mathbf{f}_{im}'$ is the $m$-th row of\/ $\mathbf{F}_i$.
Define $\tilde{\mathbf{y}}_i$, $\widetilde{\mathbf{F}}_i$, and $\tilde{\boldsymbol{\eta}}_i$
by \eqref{eq:5.2.2}. Then:
\begin{enumerate}[label=(\alph*)]
\item the fixed-effects estimator \eqref{eq:5.2.4} is consistent and asymptotically normal with
\begin{equation}
  \avar(\hat{\boldsymbol{\beta}}_{\mathrm{FE}})
  = \bigl[\E(\widetilde{\mathbf{F}}_i'\widetilde{\mathbf{F}}_i)\bigr]^{-1}\,
    \E\!\bigl[\widetilde{\mathbf{F}}_i'\,\E(\tilde{\boldsymbol{\eta}}_i\tilde{\boldsymbol{\eta}}_i')\widetilde{\mathbf{F}}_i\bigr]\,
    \bigl[\E(\widetilde{\mathbf{F}}_i'\widetilde{\mathbf{F}}_i)\bigr]^{-1}.
  \label{eq:5.2.7} \tag{5.2.7}
\end{equation}

\item This asymptotic variance is consistently estimated by
\begin{equation}
  \widehat{\avar}(\hat{\boldsymbol{\beta}}_{\mathrm{FE}})
  = \biggl(\frac{1}{n}\sum_{i=1}^{n}\widetilde{\mathbf{F}}_i'\widetilde{\mathbf{F}}_i\biggr)^{\!-1}
    \biggl(\frac{1}{n}\sum_{i=1}^{n}\widetilde{\mathbf{F}}_i'\hat{\mathbf{V}}\widetilde{\mathbf{F}}_i\biggr)
    \biggl(\frac{1}{n}\sum_{i=1}^{n}\widetilde{\mathbf{F}}_i'\widetilde{\mathbf{F}}_i\biggr)^{\!-1},
  \label{eq:5.2.8} \tag{5.2.8}
\end{equation}
\end{enumerate}
\end{proposition}
\noindent
where $\hat{\mathbf{V}}$ is the sample cross-moment matrix of transformed residuals associated
with the fixed-effects estimator:
\begin{equation}
  \underset{(M \times M)}{\hat{\mathbf{V}}}
  \equiv \frac{1}{n}\sum_{i=1}^{n}
  (\tilde{\mathbf{y}}_i - \widetilde{\mathbf{F}}_i\hat{\boldsymbol{\beta}}_{\text{FE}})
  (\tilde{\mathbf{y}}_i - \widetilde{\mathbf{F}}_i\hat{\boldsymbol{\beta}}_{\text{FE}})'.
  \label{eq:5.2.9} \tag{5.2.9}
\end{equation}

The only nontrivial part of the proof is to show:
\begin{enumerate}
\item[(1)] $\E(\widetilde{\mathbf{F}}_i'\widetilde{\mathbf{F}}_i)$ is nonsingular (and hence invertible),
\item[(2)] $\E(\mathbf{F}_i'\mathbf{Q}\boldsymbol{\eta}_i) = \mathbf{0}$ (which is needed for consistency), and
\item[(3)] $\E(\widetilde{\mathbf{F}}_i'\tilde{\boldsymbol{\eta}}_i\tilde{\boldsymbol{\eta}}_i'\widetilde{\mathbf{F}}_i)$
  (which is the Avar of $\frac{1}{n}\sum_{i=1}^{n}\widetilde{\mathbf{F}}_i'\tilde{\boldsymbol{\eta}}_i$)
  $= \E\!\bigl[\widetilde{\mathbf{F}}_i'\,\E(\tilde{\boldsymbol{\eta}}_i\tilde{\boldsymbol{\eta}}_i')\widetilde{\mathbf{F}}_i\bigr]$.
\end{enumerate}

The rest of the proof is an easy application of Kolmogorov's LLN and the Lindeberg-Levy CLT.

Proving that \eqref{eq:5.1.15} implies (1) is left as an optional analytical exercise.

To prove (2), use formula (4.6.16b) and rewrite the expectation as
\begin{align}
  \E(\mathbf{F}_i'\mathbf{Q}\boldsymbol{\eta}_i)
  &= \E\biggl(\sum_{m=1}^{M}\sum_{h=1}^{M}q_{mh}\,\mathbf{f}_{im}\cdot\eta_{ih}\biggr)
  \quad (q_{mh} \equiv (m,h)\text{ element of }\mathbf{Q})
  \notag \\
  &= \sum_{m=1}^{M}\sum_{h=1}^{M}q_{mh}\,\E(\mathbf{f}_{im}\cdot\eta_{ih}).
  \label{eq:5.2.10} \tag{5.2.10}
\end{align}
This is zero by the orthogonality conditions \eqref{eq:5.2.6}.

Proof of (3) is as follows. First note that
\[
  \E(\widetilde{\mathbf{F}}_i'\tilde{\boldsymbol{\eta}}_i\tilde{\boldsymbol{\eta}}_i'\widetilde{\mathbf{F}}_i)
  = \E\!\bigl[\widetilde{\mathbf{F}}_i'\,\E(\tilde{\boldsymbol{\eta}}_i\tilde{\boldsymbol{\eta}}_i' \mid \mathbf{x}_i)\,\widetilde{\mathbf{F}}_i\bigr].
\]
(This follows by the Law of Total Expectations and by the fact that $\mathbf{x}_i$
$[= \text{union of } (\mathbf{z}_{i1}, \ldots, \mathbf{z}_{iM})]$ includes the elements of $\mathbf{F}_i$ and
$\widetilde{\mathbf{F}}_i$ is a function of $\mathbf{F}_i$.) But
$\E(\tilde{\boldsymbol{\eta}}_i\tilde{\boldsymbol{\eta}}_i' \mid \mathbf{x}_i)$
equals the unconditional mean $\E(\tilde{\boldsymbol{\eta}}_i\tilde{\boldsymbol{\eta}}_i')$ because
\begin{align}
  \E(\tilde{\boldsymbol{\eta}}_i\tilde{\boldsymbol{\eta}}_i' \mid \mathbf{x}_i)
  &= \E(\tilde{\boldsymbol{\varepsilon}}_i\tilde{\boldsymbol{\varepsilon}}_i' \mid \mathbf{x}_i)
  \quad (\text{since } \tilde{\boldsymbol{\varepsilon}}_i
  \equiv \mathbf{Q}\boldsymbol{\varepsilon}_i
  = \mathbf{Q}(\mathbf{1}_M \cdot \alpha_i + \boldsymbol{\eta}_i)
  = \mathbf{Q}\boldsymbol{\eta}_i = \tilde{\boldsymbol{\eta}}_i)
  \notag \\
  &= \mathbf{Q}\,\E(\boldsymbol{\varepsilon}_i\boldsymbol{\varepsilon}_i' \mid \mathbf{x}_i)\,\mathbf{Q}
  \notag \\
  &= \mathbf{Q}\,\E(\boldsymbol{\varepsilon}_i\boldsymbol{\varepsilon}_i')\,\mathbf{Q}
  \quad (\text{by \eqref{eq:5.1.5}})
  \notag \\
  &= \E(\tilde{\boldsymbol{\varepsilon}}_i\tilde{\boldsymbol{\varepsilon}}_i')
  \notag \\
  &= \E(\tilde{\boldsymbol{\eta}}_i\tilde{\boldsymbol{\eta}}_i')
  \quad (\text{since } \tilde{\boldsymbol{\varepsilon}}_i = \tilde{\boldsymbol{\eta}}_i).
  \label{eq:5.2.11} \tag{5.2.11}
\end{align}

A glance at the expression for the asymptotic variance should make you suspect
that there must be some other estimator that is asymptotically more efficient,
because the asymptotic variance of an efficient estimator can be written as the
inverse of a matrix. Indeed, as you will prove in Analytical Exercise~2, there is
available an estimator that is asymptotically more efficient than the fixed-effects
estimator.

\subsection*{Digression: When $\boldsymbol{\eta}_i$ Is Spherical}

The usual error-components model assumes, additionally, that the second error
component $\boldsymbol{\eta}_i$ is spherical:
\begin{equation}
  \E(\boldsymbol{\eta}_i\boldsymbol{\eta}_i') = \sigma_\eta^2 \mathbf{I}_M,
  \quad \text{so that }
  \E(\tilde{\boldsymbol{\eta}}_i\tilde{\boldsymbol{\eta}}_i') = \sigma_\eta^2 \mathbf{Q}.
  \label{eq:5.2.12} \tag{5.2.12}
\end{equation}

If, as in many panel data sets, $m$ is time, this assumption is often referred to as the
assumption of no ``serial correlation.'' The lack of serial correlation in this sense
should not be confused with the condition that $\boldsymbol{\eta}_i$ be uncorrelated with $\boldsymbol{\eta}_j$
for $i \neq j$. The latter is a consequence of our maintained assumption that the sample is i.i.d.

Substituting this into \eqref{eq:5.2.7} and noting that
$\mathbf{Q}\widetilde{\mathbf{F}}_i = \widetilde{\mathbf{F}}_i$, we find that the expression
for the asymptotic variance simplifies to
\begin{equation}
  \avar(\hat{\boldsymbol{\beta}}_{\text{FE}})
  = \sigma_\eta^2 \cdot \bigl[\E(\widetilde{\mathbf{F}}_i'\widetilde{\mathbf{F}}_i)\bigr]^{-1}.
  \label{eq:5.2.13} \tag{5.2.13}
\end{equation}

If there is a consistent estimate, $\hat{\sigma}_\eta^2$, of $\sigma_\eta^2$, then the asymptotic variance is consistently
estimated by
\begin{equation}
  \widehat{\avar}(\hat{\boldsymbol{\beta}}_{\text{FE}})
  = \hat{\sigma}_\eta^2 \cdot
  \biggl(\frac{1}{n}\sum_{i=1}^{n}\widetilde{\mathbf{F}}_i'\widetilde{\mathbf{F}}_i\biggr)^{\!-1},
  \label{eq:5.2.14} \tag{5.2.14}
\end{equation}
which equals $n$ times $\hat{\sigma}_\eta^2 \cdot (\widetilde{\mathbf{F}}'\widetilde{\mathbf{F}})^{-1}$,
that is, $n$ times the usual expression for the estimated variance matrix when OLS is applied to the pooled sample
$(\tilde{\mathbf{y}}, \widetilde{\mathbf{F}})$ of size $Mn$.

To extend the OLS analogy, define $SSR$ as
\begin{equation}
  SSR = (\tilde{\mathbf{y}} - \widetilde{\mathbf{F}}\hat{\boldsymbol{\beta}}_{\text{FE}})'
  (\tilde{\mathbf{y}} - \widetilde{\mathbf{F}}\hat{\boldsymbol{\beta}}_{\text{FE}})
  = \sum_{i=1}^{n}(\tilde{\mathbf{y}}_i - \widetilde{\mathbf{F}}_i\hat{\boldsymbol{\beta}}_{\text{FE}})'
  (\tilde{\mathbf{y}}_i - \widetilde{\mathbf{F}}_i\hat{\boldsymbol{\beta}}_{\text{FE}}).
  \label{eq:5.2.15} \tag{5.2.15}
\end{equation}

In the pooled OLS on deviations from group means, there are $\#\boldsymbol{\beta}$ coefficients to be
estimated. The usual OLS estimate of the error variance is
\begin{equation}
  \frac{SSR}{Mn - \#\boldsymbol{\beta}}.
  \label{eq:5.2.16} \tag{5.2.16}
\end{equation}

This, however, is \emph{not} a consistent estimator of $\sigma_\eta^2$. You will be asked to show in Analytical
Exercise~3 that a consistent estimator is
\begin{equation}
  \frac{SSR}{Mn - n - \#\boldsymbol{\beta}}.
  \label{eq:5.2.17} \tag{5.2.17}
\end{equation}

(For that matter, $SSR/(Mn - n)$ is also consistent.) The reason $n$ has to be subtracted
from the denominator of \eqref{eq:5.2.16} is that $M$ transformed equations are not
linearly independent; for both sides of the $M$ transformed equations, the sum is
always zero, as you can see by multiplying both sides of \eqref{eq:5.2.1} from left by $\mathbf{1}_M'$
and making use of $\mathbf{1}_M'\mathbf{Q} = \mathbf{0}'$. The effective sample size of the pooled sample is
$Mn - n$, not $Mn$. Use of \eqref{eq:5.2.16} instead of \eqref{eq:5.2.17} is a very common mistake,
which results in an underestimation of standard errors and an overestimation of
$t$-values. For example, if $M = 3$, $n = 1{,}000$, and $\#\boldsymbol{\beta} = 5$, the $t$-values will be
overestimated by about 23\% ($=$ square root of 2995/1995 minus one).

\subsection*{Random Effects versus Fixed Effects}

Now get back to the general case of no restriction on $\E(\boldsymbol{\eta}_i\boldsymbol{\eta}_i')$, so the error term
can have an arbitrary pattern of ``serial correlation'' across $m$. Comparing \eqref{eq:5.2.6}
with \eqref{eq:5.1.8}, we see that the orthogonality conditions \emph{not} used by the fixed-effects
estimator are
\begin{equation}
  \E(\mathbf{f}_{im} \cdot \alpha_i) = \mathbf{0} \text{ for all } m, \quad
  \E(\mathbf{b}_i \cdot \alpha_i) = \mathbf{0}, \quad
  \E(\mathbf{b}_i \cdot \eta_{im}) = \mathbf{0} \text{ for all } m.
  \label{eq:5.2.18} \tag{5.2.18}
\end{equation}

That is, the fixed-effects estimator is robust to the failure of \eqref{eq:5.2.18}.
Let $\hat{\boldsymbol{\beta}}_{\text{RE}}$ be the elements of
$\hat{\boldsymbol{\delta}}_{\text{RE}}$ (the random-effects estimator applied to the original $M$-equation
system \eqref{eq:5.1.1pp} on page~329) that correspond to $\boldsymbol{\beta}$:
\begin{equation}
  \hat{\boldsymbol{\delta}}_{\text{RE}} \equiv
  \begin{bmatrix} \hat{\boldsymbol{\beta}}_{\text{RE}} \\ \hat{\boldsymbol{\gamma}}_{\text{RE}} \end{bmatrix}.
  \label{eq:5.2.19} \tag{5.2.19}
\end{equation}

Also, let $\avar(\hat{\boldsymbol{\beta}}_{\text{RE}})$ be the asymptotic variance of
$\hat{\boldsymbol{\beta}}_{\text{RE}}$. It is the leading submatrix
of $\avar(\hat{\boldsymbol{\delta}}_{\text{RE}})$ (given by (4.6.9$'$) on page~293)
corresponding to $\boldsymbol{\beta}$.

The random-effects estimator $\hat{\boldsymbol{\beta}}_{\text{RE}}$ is an efficient estimator while
$\hat{\boldsymbol{\beta}}_{\text{FE}}$ is consistent
but not efficient. However, if \eqref{eq:5.2.18} is violated,
$\hat{\boldsymbol{\beta}}_{\text{RE}}$ is no longer guaranteed
to be consistent, while $\hat{\boldsymbol{\beta}}_{\text{FE}}$ remains consistent. Thus, a natural test of \eqref{eq:5.2.18} is to
consider the difference between the two estimators,
\begin{equation}
  \hat{\mathbf{q}} \equiv \hat{\boldsymbol{\beta}}_{\text{FE}} - \hat{\boldsymbol{\beta}}_{\text{RE}}.
  \label{eq:5.2.20} \tag{5.2.20}
\end{equation}

It is easy to prove that $\sqrt{n}\,\hat{\mathbf{q}}$ is asymptotically normal. The Hausman principle
(originally developed for ML estimators but proved to be applicable to GMM estimators
in Analytical Exercise~9 to Chapter~3) implies that
\begin{equation}
  \avar(\hat{\mathbf{q}})
  = \avar(\hat{\boldsymbol{\beta}}_{\text{FE}}) - \avar(\hat{\boldsymbol{\beta}}_{\text{RE}}).
  \label{eq:5.2.21} \tag{5.2.21}
\end{equation}

(So, there is no need to incorporate the asymptotic covariance between
$\hat{\boldsymbol{\beta}}_{\text{FE}}$ and
$\hat{\boldsymbol{\beta}}_{\text{RE}}$ because it is zero.) By Proposition~4.7, the leading submatrix of (4.6.9$'$)
on page~293 provides a consistent estimator of $\avar(\hat{\boldsymbol{\beta}}_{\text{RE}})$.
By Proposition~5.1,
$\avar(\hat{\boldsymbol{\beta}}_{\text{FE}})$ is consistently estimated by \eqref{eq:5.2.8}.
A consistent estimator of $\avar(\hat{\mathbf{q}})$,
therefore, is the difference. Write this as $\widehat{\avar}(\hat{\mathbf{q}})$.
It is proved in the appendix
that (i) $\avar(\hat{\mathbf{q}})$ is nonsingular (and hence positive definite) and (ii) the Hausman
statistic defined below is guaranteed to be nonnegative for any sample $\{\mathbf{y}_i, \mathbf{Z}_i\}$. To
summarize,

\begin{proposition}[Hausman specification test]
\label{prop:5.2}
Suppose the assumptions of the
error-components model (\eqref{eq:5.1.1pp} on page~329, \eqref{eq:5.1.2}, \eqref{eq:5.1.8},
\eqref{eq:5.1.4}--\eqref{eq:5.1.6},
and \eqref{eq:5.1.15}) hold. Define $\hat{\mathbf{q}}$ and $\widehat{\avar}(\hat{\mathbf{q}})$
as just described. Then the Hausman statistic
\begin{equation}
  H \equiv n \cdot \hat{\mathbf{q}}'\bigl(\widehat{\avar}(\hat{\mathbf{q}})\bigr)^{-1}\hat{\mathbf{q}}
  \label{eq:5.2.22} \tag{5.2.22}
\end{equation}
is distributed asymptotically chi-square with $\#\boldsymbol{\beta}$ degrees of freedom.
It is nonnegative in finite samples.
\end{proposition}

This test is a specification test because it can detect a failure of \eqref{eq:5.2.18} which is
part of the maintained assumptions of the error-components model. More specifically,
consider a sequence of local alternatives that satisfy all the assumptions
of the error-components model except \eqref{eq:5.2.18}.\footnote{The notion of local alternatives
was introduced in Section~2.4.}
With some additional technical
assumptions, it can be shown that the Hausman statistic converges to a noncentral
$\chi^2$ distribution with a noncentrality parameter along those sequences of local alternatives.\footnote{See Newey (1985). The required technical assumption is Newey's Assumption~5.}
That is, the Hausman statistic has power against local alternatives under
which the random-effects estimator is inconsistent.

\subsection*{Relaxing Conditional Homoskedasticity}

It is straightforward to derive the large sample distribution of the fixed-effects estimator
without conditional homoskedasticity. As will be seen in the next section,
extension of the large sample results to cover unbalanced panels is actually easier
without conditional homoskedasticity.

\begin{proposition}[fixed-effects estimator without conditional homoskedasticity]
\label{prop:5.3}
Drop conditional homoskedasticity \eqref{eq:5.1.5} from the hypothesis of Proposition~5.1.
Define $\tilde{\mathbf{y}}_i$, $\widetilde{\mathbf{F}}_i$, and $\tilde{\boldsymbol{\eta}}_i$
by \eqref{eq:5.2.2}. Then:
\begin{enumerate}[label=(\alph*)]
\item the fixed-effects estimator \eqref{eq:5.2.4} is consistent and asymptotically normal with
\begin{equation}
  \avar(\hat{\boldsymbol{\beta}}_{\mathrm{FE}})
  = \bigl[\E(\widetilde{\mathbf{F}}_i'\widetilde{\mathbf{F}}_i)\bigr]^{-1}\,
    \E(\widetilde{\mathbf{F}}_i'\tilde{\boldsymbol{\eta}}_i\tilde{\boldsymbol{\eta}}_i'\widetilde{\mathbf{F}}_i)\,
    \bigl[\E(\widetilde{\mathbf{F}}_i'\widetilde{\mathbf{F}}_i)\bigr]^{-1}.
  \label{eq:5.2.23} \tag{5.2.23}
\end{equation}

\item If, furthermore, some appropriate finite fourth-moment assumption (which is
  an adaptation of Assumption~4.6) is satisfied, then the asymptotic variance is
  consistently estimated by
\begin{equation}
  \widehat{\avar}(\hat{\boldsymbol{\beta}}_{\mathrm{FE}})
  = \biggl(\frac{1}{n}\sum_{i=1}^{n}\widetilde{\mathbf{F}}_i'\widetilde{\mathbf{F}}_i\biggr)^{\!-1}
    \biggl(\frac{1}{n}\sum_{i=1}^{n}\widetilde{\mathbf{F}}_i'\hat{\boldsymbol{\eta}}_i\hat{\boldsymbol{\eta}}_i'\widetilde{\mathbf{F}}_i\biggr)
    \biggl(\frac{1}{n}\sum_{i=1}^{n}\widetilde{\mathbf{F}}_i'\widetilde{\mathbf{F}}_i\biggr)^{\!-1},
  \label{eq:5.2.24} \tag{5.2.24}
\end{equation}
where $\hat{\boldsymbol{\eta}}_i \equiv \tilde{\mathbf{y}}_i - \widetilde{\mathbf{F}}_i\hat{\boldsymbol{\beta}}_{\mathrm{FE}}$
(this $\hat{\boldsymbol{\eta}}_i$ should not be confused with
$\tilde{\boldsymbol{\eta}}_i$ in \eqref{eq:5.2.1}).
\end{enumerate}
\end{proposition}

Again, there are two ways to prove this. Write the fixed-effects estimator as a
GMM estimator and then apply the large sample theory of GMM estimators. Or
you can prove it directly by inspecting the expression \eqref{eq:5.2.4} for the sampling error.
The only non-trivial part of the proof is that the middle summation in the expression
for $\widehat{\avar}(\hat{\boldsymbol{\beta}}_{\text{FE}})$ converges in probability to its population counterpart,
$\E(\widetilde{\mathbf{F}}_i'\tilde{\boldsymbol{\eta}}_i\tilde{\boldsymbol{\eta}}_i'\widetilde{\mathbf{F}}_i)$.
The required technique is very similar to the one used for proving Proposition~4.2
and so is not repeated here.

\bigskip
\noindent\rule{\textwidth}{0.5pt}
\textsc{Questions for Review}
\medskip

\begin{enumerate}
\item For the large sample results of this section to hold, do we have to assume that
  $\alpha_i$ and $\boldsymbol{\eta}_i$ are orthogonal to each other?

\item (Dropping redundant equations)\quad As mentioned in the text, the $M$ transformed
  equations are not linearly independent in the sense that any one of the $M$ transformed
  equations is a linear combination of the remaining $M-1$ equations. So
  consider forming a pooled sample of size $Mn - n$ by dropping one transformed
  equation (say, the last equation) for each group ($i$). Is the OLS estimator on
  this smaller pooled sample the fixed-effects estimator? [Answer: No.]

\item (Importance of cross orthogonalities)\quad Suppose \eqref{eq:5.2.6} holds for $m = h$ but
  not necessarily for $m \neq h$. Would the fixed-effects estimator be necessarily
  consistent? [Answer: No.]

\item Verify that $\E(\tilde{\boldsymbol{\eta}}_i\tilde{\boldsymbol{\eta}}_i')$ is singular.
  (Nevertheless, $\E\!\bigl[\widetilde{\mathbf{F}}_i'\,\E(\tilde{\boldsymbol{\eta}}_i\tilde{\boldsymbol{\eta}}_i')\widetilde{\mathbf{F}}_i\bigr]$
  in \eqref{eq:5.2.7} and
  $\E(\widetilde{\mathbf{F}}_i'\tilde{\boldsymbol{\eta}}_i\tilde{\boldsymbol{\eta}}_i'\widetilde{\mathbf{F}}_i)$
  in \eqref{eq:5.2.23} are nonsingular, as an optional analytical exercise will
  ask you to show.)

\item (Consistent estimation of error covariances)\quad To prove that \eqref{eq:5.2.9} is consistent
  for $\E(\tilde{\boldsymbol{\eta}}_i\tilde{\boldsymbol{\eta}}_i')$, you would apply Proposition~4.1 to the transformed system of
  $M$ equations. Verify that the conditions of the proposition are satisfied here.
  In particular, is the requirement that $\E(\mathbf{f}_{im}\tilde{\mathbf{f}}_{im}')$
  (where $\tilde{\mathbf{f}}_{im}$ is the $m$-th row of
  $\widetilde{\mathbf{F}}_i$) be finite satisfied by the error-component model?
  \textbf{Hint:} $\tilde{\mathbf{f}}_{im}$ is a linear
  transformation of $\mathbf{x}_i$. \eqref{eq:5.1.5} and \eqref{eq:5.1.6}
  imply that $\E(\mathbf{x}_i\mathbf{x}_i')$ is nonsingular.

\item (What the Hausman statistic tests)\quad For simplicity, assume
  $\boldsymbol{\Sigma}\;(\equiv \E(\boldsymbol{\varepsilon}_i\boldsymbol{\varepsilon}_i'))$ is
  known. Also for simplicity, assume that there are no common regressors, so
  $\mathbf{z}_{im} = \mathbf{f}_{im}$ and \eqref{eq:5.2.18} reduces to
  $\E(\mathbf{z}_{im} \cdot \alpha_i) = \mathbf{0}$ for all $m$. This is the
  restriction \emph{not} used by the fixed-effects estimator. Is the random-effects estimator
  necessarily inconsistent when \eqref{eq:5.1.8a} fails but all the other assumptions
  of the error-components model (such as \eqref{eq:5.1.8b}) are satisfied?
  \textbf{Hint:} Show
  that, without \eqref{eq:5.1.8a},
  \[
    \plim \hat{\boldsymbol{\delta}}_{\text{RE}} - \boldsymbol{\delta}
    = (\mathbf{Z}_i'\boldsymbol{\Sigma}^{-1}\mathbf{Z}_i)^{-1}\,
      \E(\mathbf{Z}_i'\boldsymbol{\Sigma}^{-1}\boldsymbol{\varepsilon}_i),
  \]
  where $\boldsymbol{\varepsilon}_i = \mathbf{1}_M \cdot \alpha_i + \boldsymbol{\eta}_i$.
  Provided \eqref{eq:5.1.8b} holds, \eqref{eq:5.1.8a} is \emph{sufficient} for
  $\E(\mathbf{Z}_i'\boldsymbol{\Sigma}^{-1}\boldsymbol{\varepsilon}_i)$ to be zero.
  Is it necessary? Thus, strictly speaking, the Hausman
  test is not a test of \eqref{eq:5.1.8a}; it is a test of
  $\E(\mathbf{Z}_i'\boldsymbol{\Sigma}^{-1}\boldsymbol{\varepsilon}_i) = \mathbf{0}$.
\end{enumerate}
